\pendiente{OBLIGATORIO revisar con el tutor. La Normativa de TFG de la EPS (8/5/2026,
art. 5.1.3) exige declarar, al solicitar la defensa, que se han respetado las condiciones de
uso de herramientas de IA especificadas por el ED-TFG \cite{eps_normativa_tfg_2026}. Esas
condiciones se publican en la plataforma de TFG, que requiere autenticación y no se ha
podido consultar: el autor debe contrastar esta sección con ellas y adaptarla.}

Este trabajo se ha desarrollado con herramientas de \acs{ia} generativa de forma central, no
accesoria, y así se declara. Esta sección resume qué se hizo con ellas, bajo qué control y
cómo se verificó el resultado; el detalle, con las instrucciones empleadas y los errores
registrados, está en el \cref{AP:AI}.

\subsection{Herramientas}

\begin{itemize}
  \item \textbf{Claude Code} (Anthropic) \cite{claude_code_docs}, con modelos de la familia
        Claude Opus, como \dfn{agente} de la línea del núcleo docente, responsable de la
        base común y editor de esta memoria.
  \item \textbf{Codex} (OpenAI) \cite{openai_codex}, configurado con el modelo
        \texttt{gpt-6.1-sol}, como \dfn{agente} de la línea de gestión, redactor de
        borradores y responsable de publicar la demostración.
\end{itemize}

\subsection{Qué se automatizó}

Los agentes ejecutaron la mayor parte del trabajo material: la lectura del repositorio y del
contrato del \dfn{backend}, la escritura del código de la aplicación y de sus pruebas, los
\textit{scripts} de análisis y de auditoría, la documentación técnica y los registros de
coordinación, y los borradores de los capítulos de esta memoria. El estado del arte se
redactó a partir de borradores de Codex revisados por la línea del núcleo; el resto de
capítulos los redactó Claude Code, salvo la parte de la implementación y de las pruebas
relativa a la línea de gestión, que parte de borradores de Codex. Los \textit{commits} se crearon con la identidad Git del
autor y, por una regla del repositorio, sus mensajes no mencionan las herramientas: esta
sección y el \cref{AP:AI} son, por tanto, la declaración de su uso, y el reparto por líneas
de \cref{TAB:IMP-SPLIT} permite atribuir cada parte del código a su agente.

\subsection{Qué decidió el autor}

El autor definió los objetivos y el alcance, eligió las herramientas y la forma de trabajo,
fijó las reglas que los agentes debían cumplir (no inventar operaciones de la \acs{api}, no
modificar rutas ajenas, no publicar ni reescribir historial sin autorización), tomó las
decisiones de producto y aceptó o rechazó los resultados. El rechazo de la primera dirección
visual, que obligó a rehacer el sistema de diseño (\cref{SEC:DES-VISUAL}), es el ejemplo más
claro de una decisión que ninguna comprobación automática habría tomado. La autorización
para publicar la demostración también fue suya.

\subsection{Cómo se verificó}

El principio fue no aceptar una afirmación de un agente que no pudiera comprobarse en el
repositorio. En la práctica:
\begin{itemize}
  \item cada iteración debía superar las mismas comprobaciones deterministas (análisis
        estático, pruebas, compilación, auditorías de estilo y de propiedad) antes de
        consolidarse (\cref{SEC:TEST-AUDITS});
  \item las operaciones de la \acs{api} se contrastaron con el documento OpenAPI real y, en
        los casos dudosos, con el código de las vistas de Django (\cref{SEC:AN-BACKEND});
  \item cada línea revisó el trabajo de la otra: la del núcleo auditaba la rama de gestión, y
        la de gestión detectó carencias de la base común que la del núcleo corrigió
        (\cref{SEC:IMP-PARALLEL});
  \item la redacción de la memoria se contrastó con el código y con el historial, lo que
        sirvió también para detectar defectos (el autoguardado de anotaciones,
        \cref{SEC:IMP-CRITICAL}) y datos erróneos en la propia documentación, como un
        recuento de operaciones del contrato que se corrigió de 96 a 95.
\end{itemize}

\subsection{Errores atribuibles a los agentes}

El uso de agentes no eliminó los errores, y declararlos forma parte de la honestidad que
exige la normativa. Los registrados incluyen: un \textit{commit} consolidado con pruebas en
rojo por no comprobar el código de salida de la verificación; una dirección visual que
copiaba la paleta de la referencia en lugar de su principio de diseño; los dos defectos de
autoguardado descritos en \cref{SEC:IMP-CRITICAL}; una regresión del editor de permisos en la
línea de gestión; y un párrafo de un borrador del estado del arte que describía carencias de
la base común ya corregidas. Todos se detectaron mediante pruebas, auditorías, la revisión
cruzada o la revisión del autor, y se corrigieron con su prueba de regresión cuando
procedía (\cref{AP:ITERATIONS}).

\subsection{Límites}

Los textos de esta memoria son borradores avanzados que el autor debe revisar, corregir y
asumir como propios; las marcas de pendiente señalan lo que solo él puede aportar. Ninguna
afirmación sobre resultados se apoya en mediciones que no consten en el repositorio: no se
han hecho pruebas con usuarios ni, a la fecha de redacción, una validación completa contra el
\dfn{backend} real (\cref{SEC:TEST-REAL}).

\pendiente{El autor debe completar con su propia valoración: grado de revisión que aplicó a
cada parte (código, pruebas, memoria), qué aprendió del proceso y cómo encaja este uso con
las condiciones del ED-TFG.}
